\documentclass{article}
\usepackage[utf8]{inputenc}
\usepackage{amsmath}
\usepackage{geometry}
\geometry{margin=2cm}
\begin{document}

\section*{Resolução da questão: impacto da droga citocalasina B no citosqueleto celular}

	extbf{Interpretação + Estratégia + Análise}

A questão avalia o entendimento do impacto de drogas específicas nos componentes do citosqueleto celular, principalmente o efeito da citocalasina B em células de ameba. Ao observar a imagem fornecida, percebe-se que as células tratadas apresentaram alteração no formato e na capacidade de movimento.

Drogas como a citocalasina B atuam na despolimerização dos microfilamentos de actina, que são fi lamentos proteicos essenciais para manter a forma celular e possibilitar o movimento. Assim, relembrando a ação específica da citocalasina B, relaciona-se sua atuação com a perda de integridade dessas estruturas internas.

O passo estratégico consiste em associar a mudança visual na célula à ação da droga, que impede a formação adequada e a manutenção dos microfilamentos.

\textbf{Conteúdo + Mini revisão}

O citosqueleto é uma rede de filamentos internos que fornece suporte, forma e movimentação às células. Os principais componentes do citosqueleto incluem os microtúbulos, microfilamentos de actina e filamentos intermediários.

Os microfilamentos, compostos por actina, são cruciais na sustentação da forma da célula e na movimentação celular. Drogas específicas, como a citocalasina B, atuam na decomposição ou desorganização destes microfilamentos, causando alterações na morfologia e no movimento.

\textbf{Resolução + Pulo do Gato}

Analisando a imagem, a célula tratada mostrou significativa perda de forma e movimento, indicando que o suporte estrutural foi afetado. Como a citocalasina B atua na despolimerização dos microfilamentos de actina, podemos concluir que este componente foi o alvo da droga.

\textbf{Armado com conceitos}

A alternativa correta é a letra C: Microfilamentos.

\textbf{Código da questão}

CIEN	extunderscore CEL	extunderscore 001

\end{document}